\section{Same finite nonflat algebra principal degree over the original curve ring}

\begin{definition}[The actual whole scalar module quotient is the original ring quotient]
\label{def:fc-finite-algebra-scalar-quotient}
\leanok
\lean{EnvelopeProperty.FiniteAlgebraPrincipalDegrees.scalarQuotientIso}
For an original algebra and original scalar, the quotient by the
original scalar endomorphism is linearly isomorphic over the
original base ring to the SAME original algebra quotient by the
principal ideal of that scalar's actual image.
\end{definition}
\begin{proof}
The original image of scalar multiplication is exactly the
restriction of the SAME original principal algebra ideal. The
actual quotient-by-equal-submodules and original restriction of
scalars quotient isomorphisms give the entire original quotient
comparison.
\end{proof}

\begin{theorem}[Same actual finite algebra principal degree, without flatness]
\label{thm:fc-finite-algebra-principal-degree}
\leanok
\lean{EnvelopeProperty.FiniteAlgebraPrincipalDegrees.length_principalQuotient_eq_rank_mul_ord, EnvelopeProperty.FiniteAlgebraPrincipalDegrees.length_principalQuotient_eq_fieldDegree_mul_ord}
\uses{def:fc-finite-algebra-scalar-quotient,thm:fc-principal-module-colength-generic-rank}
For an actual finite injective extension of original domains, with
the original base Noetherian of dimension at most one, the original
base length of the SAME algebra quotient by any nonzero original
scalar is its actual generic rank times the native original ring
order of that SAME scalar. With original compatible fraction fields,
the rank is the degree of that SAME original fraction-field
extension. No normality, flatness, basis or length certificate is
assumed.
\end{theorem}
\begin{proof}
Actual original injectivity and the domain structure derive original
torsion-freeness of the entire original algebra module. The genuine
principal-module colength theorem internally produces a full-rank
free lattice of finite colength and derives the length formula.
The preceding whole quotient comparison retains the SAME actual
principal algebra quotient. Native fraction-field rank comparison,
with the literal original coefficient towers, identifies the
multiplier with the original field extension degree.
\end{proof}
